% =============================================================================
% COPYRIGHT AND DISCLAIMER
% =============================================================================
% This file is provided as a convenience template for competition participants.
% By using it, you accept full responsibility for compliance with all applicable
% copyright, licensing, and usage rules of the Association for Computing
% Machinery (ACM) and of any LaTeX packages you load (including, without
% limitation, the "acmart" document class and related ACM materials).
%
% Consult ACM's official publication, template, and copyright policies before
% using this setup for any purpose other than this competition. Do not assume
% that this template may be reused, redistributed, or submitted elsewhere
% without independently verifying ACM's current rules.
%
% Use of the "acmart" class, the "sigconf" option, or ACM-style formatting in
% this template does NOT mean that this report, this competition, its organizers,
% or any resulting PDF is sponsored, endorsed, affiliated with, or published by
% ACM. ACM trademarks and branding remain the property of ACM.
%
% Template packages live in preamble.tex.
% Citations live in main.bib. Edit this file with your paper content.
% =============================================================================

\documentclass[sigconf]{acmart}
% =============================================================================
% Template setup — participants normally do not edit this file.
% Author content belongs in main.tex; citations belong in main.bib.
% \documentclass lives in main.tex so Overleaf does not treat this file as root.
% =============================================================================

% --- Competition / non-ACM-submission friendly rights metadata -------------
\setcopyright{none}
\settopmatter{printacmref=false}
\acmConference[Kaggle Competition]{Filament Segmentation Challenge}{2026}{\href{https://www.kaggle.com/competitions/filament-segmentation-2026}{kaggle.com/competitions/filament-segmentation-2026}}
\acmBooktitle{Competition Report}
\acmDOI{}
\acmISBN{}

% --- Author guidance text --------------------------------------------------
% \guide{...} is gray italic while \guideStyletrue; omitted from the PDF if false.
% Flip the switch in main.tex after % =============================================================================
% Template setup — participants normally do not edit this file.
% Author content belongs in main.tex; citations belong in main.bib.
% \documentclass lives in main.tex so Overleaf does not treat this file as root.
% =============================================================================

% --- Competition / non-ACM-submission friendly rights metadata -------------
\setcopyright{none}
\settopmatter{printacmref=false}
\acmConference[Kaggle Competition]{Filament Segmentation Challenge}{2026}{\href{https://www.kaggle.com/competitions/filament-segmentation-2026}{kaggle.com/competitions/filament-segmentation-2026}}
\acmBooktitle{Competition Report}
\acmDOI{}
\acmISBN{}

% --- Author guidance text --------------------------------------------------
% \guide{...} is gray italic while \guideStyletrue; omitted from the PDF if false.
% Flip the switch in main.tex after % =============================================================================
% Template setup — participants normally do not edit this file.
% Author content belongs in main.tex; citations belong in main.bib.
% \documentclass lives in main.tex so Overleaf does not treat this file as root.
% =============================================================================

% --- Competition / non-ACM-submission friendly rights metadata -------------
\setcopyright{none}
\settopmatter{printacmref=false}
\acmConference[Kaggle Competition]{Filament Segmentation Challenge}{2026}{\href{https://www.kaggle.com/competitions/filament-segmentation-2026}{kaggle.com/competitions/filament-segmentation-2026}}
\acmBooktitle{Competition Report}
\acmDOI{}
\acmISBN{}

% --- Author guidance text --------------------------------------------------
% \guide{...} is gray italic while \guideStyletrue; omitted from the PDF if false.
% Flip the switch in main.tex after \input{preamble}.
%
% Do not use \guide{...} for \email{...}.
\usepackage{xcolor}
\newif\ifguideStyle
\guideStyletrue
\newcommand{\guide}[1]{%
  \ifguideStyle
    {\color{gray!65}\itshape #1}%
  \fi
}

% --- Math ------------------------------------------------------------------
% Do NOT load amssymb (or amsfonts): acmart already loads newtxmath, which
% defines the same symbols (e.g. \Bbbk) and causes a conflict on Overleaf.
% amsmath is already provided by acmart; mathtools extends it.
\usepackage{mathtools}

% --- Figures ---------------------------------------------------------------
% graphicx is already loaded by acmart
\usepackage{subcaption}

% --- Tables ----------------------------------------------------------------
% booktabs is already loaded by acmart
\usepackage{tabularx}
\usepackage{multirow}

% --- Bibliography style is set in main.tex (ACM-Reference-Format) ----------
.
%
% Do not use \guide{...} for \email{...}.
\usepackage{xcolor}
\newif\ifguideStyle
\guideStyletrue
\newcommand{\guide}[1]{%
  \ifguideStyle
    {\color{gray!65}\itshape #1}%
  \fi
}

% --- Math ------------------------------------------------------------------
% Do NOT load amssymb (or amsfonts): acmart already loads newtxmath, which
% defines the same symbols (e.g. \Bbbk) and causes a conflict on Overleaf.
% amsmath is already provided by acmart; mathtools extends it.
\usepackage{mathtools}

% --- Figures ---------------------------------------------------------------
% graphicx is already loaded by acmart
\usepackage{subcaption}

% --- Tables ----------------------------------------------------------------
% booktabs is already loaded by acmart
\usepackage{tabularx}
\usepackage{multirow}

% --- Bibliography style is set in main.tex (ACM-Reference-Format) ----------
.
%
% Do not use \guide{...} for \email{...}.
\usepackage{xcolor}
\newif\ifguideStyle
\guideStyletrue
\newcommand{\guide}[1]{%
  \ifguideStyle
    {\color{gray!65}\itshape #1}%
  \fi
}

% --- Math ------------------------------------------------------------------
% Do NOT load amssymb (or amsfonts): acmart already loads newtxmath, which
% defines the same symbols (e.g. \Bbbk) and causes a conflict on Overleaf.
% amsmath is already provided by acmart; mathtools extends it.
\usepackage{mathtools}

% --- Figures ---------------------------------------------------------------
% graphicx is already loaded by acmart
\usepackage{subcaption}

% --- Tables ----------------------------------------------------------------
% booktabs is already loaded by acmart
\usepackage{tabularx}
\usepackage{multirow}

% --- Bibliography style is set in main.tex (ACM-Reference-Format) ----------

\usepackage{tikz}  % added: pipeline diagram
\usetikzlibrary{positioning,arrows.meta}

% Guide styling: \guideStyletrue = show gray italic hints; \guideStylefalse = omit them
\guideStylefalse
% \guideStylefalse

\begin{document}

\title{Consensus U-Nets, a Seed Refiner and a Learned Filter for Solar Filament Segmentation}  % Replace this with your title
\subtitle{A Solution to the Solar Filament Segmentation Challenge 2026} % DO NOT CHANGE


\author{Nhat Vo} 
\affiliation{%
  \institution{Sensen.ai}  % <- Optional (remove if you are not currently affiliated with any institute)
  \city{Melbourne}
  \country{Australia}
}
\email{mr.nhatvo@gmail.com}   % <- Replace with your own email 

% \author{\guide{Second Author}}
% \affiliation{%
%   \institution{\guide{Your Institution}}  % <- Optional (remove if you are not currently affiliated with any institute)
%   \city{\guide{City}}
%   \country{\guide{Country}}
% }
% \email{second.author@example.com}  % <- Replace with your own email

\begin{abstract}
    This report describes our solution to the Solar Filament Segmentation Challenge~2026, a Kaggle competition on automatic segmentation of solar filaments in GONG H-$\alpha$ observations. The training labels are noisy: annotators of the same image agree with each other at only
PQ 0.34. Ten U-Nets (five folds at 1024 and five at 2048 pixels) learn the per-pixel mean of
all annotators' masks; a simple post-processing step turns their mean map into candidate
filaments; a seed refiner re-segments a crop around each candidate; and a learned filter
removes candidates that annotators are unlikely to mark. Out of fold, the pipeline reaches
PQ 0.445 on all 707 training images (95\% bootstrap interval 0.435--0.454), against 0.413 for
a single U-Net. The public leaderboard score is 0.38. The code, the weights and a notebook
that rebuilds the submission are in our repository.
\end{abstract}

\maketitle

\ifguideStyle
\noindent{\color{gray!65}\itshape
Guide mode is on.}
\par\medskip
\fi

\guide{Please limit your report to four (4) pages, excluding the Acknowledgment and References sections. Throughout this document, please follow these considerations: 
\begin{itemize}
    \item Do not alter any text that is currently in black color (except for the emails under authors' name).
    \item Replace all gray text with your own.
    \item When you are done with the report, please replace \texttt{\textbackslash guideStyletrue} with \texttt{\textbackslash guideStylefalse}, to remove any gray text. Alternatively, you can remove any text wrapped in \texttt{\textbackslash guide\{\}},
    \item Be as succinct as possible, and only provide details regarding your own work. For example, there is no need to discuss what solar filaments are, or why filament segmentation matters.
    \item If you must, you can add extra sections.
\end{itemize}}


% ==================================== %
%              INTRODUCTION            %
% ==================================== %
\section{Introduction}
    This report goes over the details of our submitted solution to the Solar Filament Segmentation Challenge~2026 \cite{filament-segmentation-2026}. The challenge, launched on July~10,~2026, is organized by the Earth-Space AI Research (ESAIR) Lab\footnote{\href{https://www.esairlab.com/home}{www.esairlab.com/}}. This challenge utilizes the MAGFiLO dataset for evaluation of the segmentation algorithms \cite{ahmadzadeh2024dataset}. 

Each test image is a $2048\times2048$ grayscale H$\alpha$ full-disk image, and a submission
gives one run-length mask per filament. The metric is Panoptic Quality
(PQ)~\cite{kirillov2019panoptic}: a prediction and a ground-truth (GT) filament match when
their IoU exceeds 0.5, and
\[
\mathrm{PQ}=\frac{\sum_{\mathrm{TP}}\mathrm{IoU}}{\mathrm{TP}+\tfrac12\mathrm{FP}+\tfrac12\mathrm{FN}}.
\]
Each annotator's labelling of an image (an \emph{entry}) is scored separately against the same
predictions, and the counts are pooled. Our main finding is that the labels limit the score:
when we score each annotator against the other annotators of the same image, PQ is only 0.340
(SQ 0.634, RQ 0.536). Our design follows from this: learn the consensus, combine two
resolutions, and reject candidates that annotators do not mark.


% ==================================== %
%              METHODOLOGY             %
% ==================================== %
\section{Methodology}\label{sec:methodology}

\subsection{Data, sampling and splits}

The training set has 707 images and 8,199 filament polygons from 37 annotators; 411 images
have one annotation set, 145 have two and 151 have three (1,154 entries). Filaments are small:
the median polygon covers 1,228 pixels (0.03\% of the image). We use five folds grouped by
image, so all entries of an image stay in one fold. Filaments last several days, so we checked
whether random folds are too easy: 28\% of validation images have a training image within one
day, but for the test images the share is 36\%. Random folds therefore match the test
situation, and folds by week would be harsher than the test. All numbers below are out of fold.

During training, 70\% of crops are centred on a random filament pixel and 30\% are placed
uniformly. Augmentations are flips, 90-degree rotations and small brightness and contrast
changes. We ported the organisers' PQ to run-length masks; a unit test checks that it matches
their notebook to $10^{-6}$.

\subsection{Pipeline}

\begin{figure}[t]
\centering
\begin{tikzpicture}[node distance=3.2mm, font=\footnotesize,
  box/.style={draw, rounded corners=1pt, align=center, inner sep=2.5pt, minimum width=36mm},
  arr/.style={-{Stealth[length=1.6mm]}}]
\node[box] (img) {H$\alpha$ image, $2048\times2048$};
\node[box, below left=4mm and -14mm of img, minimum width=30mm] (a) {5 U-Nets at 1024\\(ResNet34, folds 0--4)};
\node[box, below right=4mm and -14mm of img, minimum width=30mm] (b) {5 U-Nets at 2048\\(ResNet34, folds 0--4)};
\node[box, below=14mm of img] (mean) {mean probability map\\(10 models, h-flip TTA)};
\node[box, below=of mean] (cand) {threshold 0.5, connected regions,\\join gaps $<15$\,px, drop $<300$\,px};
\node[box, below=of cand] (ref) {seed refiner on a crop per candidate\\(image, seed, probability; 5 folds)};
\node[box, below=of ref] (filt) {learned filter: keep if predicted share\\of matching annotators $\geq 0.35$};
\node[box, below=of filt] (out) {one mask per filament, no overlaps};
\draw[arr] (img) -- (a); \draw[arr] (img) -- (b);
\draw[arr] (a) -- (mean); \draw[arr] (b) -- (mean);
\draw[arr] (mean) -- (cand); \draw[arr] (cand) -- (ref); \draw[arr] (ref) -- (filt); \draw[arr] (filt) -- (out);
\end{tikzpicture}
\caption{The end-to-end inference pipeline. Only the pixels of the test image are used.}
\label{fig:pipeline}
\end{figure}

Fig.~\ref{fig:pipeline} shows the pipeline. It uses existing building blocks: U-Net
\cite{ronneberger2015unet} with a ResNet34 encoder \cite{he2016resnet} (ImageNet start,
segmentation-models-pytorch \cite{iakubovskii2019smp}). The seed refiner is our own design.

\textbf{First stage.} Each U-Net has one input channel. Its target is the fraction of
annotators that marked each pixel, a soft label in $[0,1]$; the loss is binary cross-entropy
plus soft Dice. Per fold we train a \emph{1024 model} (image downscaled to 1024, random
$512^2$ crops, batch 8, 60 epochs) and a \emph{2048 model} (full resolution, $1024^2$ crops,
batch 4, 60 epochs). Both crops cover the same part of the disk; the 2048 model sees four
times the detail. At inference we average the logits of the image and its horizontal flip,
then the probability maps of all ten models.

\textbf{Candidates (post-processing).} The mean map is thresholded at 0.5. Connected regions
closer than 15 pixels form one candidate, which joins fragments without changing any mask
pixel; candidates under 300 pixels are removed. These three values were chosen on folds 0
and 1 from 108 settings and then fixed.

\textbf{Seed refiner.} Most first-stage errors are candidates on unmarked places or with too
little overlap, not merges (Section~\ref{sec:weak}). The refiner looks at one candidate at a
time: a square crop (1.5 times its bounding box, at least 128 pixels, resized to 256) with
three channels, namely the image, the candidate mask (the \emph{seed}) and the first-stage
probability. Its target is, for each annotator, the union of the GT filaments that touch the
candidate, averaged over the annotators; a candidate that touches no GT filament has an empty
target, so the refiner learns to reject it. The refiner is a three-channel ResNet34 U-Net,
trained for 30 epochs on the candidates of the other four folds, which come from out-of-fold
maps and so look like test candidates. At inference the five fold refiners are averaged; each
contested pixel goes to the candidate with the higher probability, so masks never overlap.

\textbf{Learned candidate filter.} The refined pipeline still has about five false positives
per image, many on places that no annotator marked. A LightGBM regressor predicts, for each
refined candidate, the share of the image's annotators whose filament it matches at IoU above
0.5, which is what PQ counts. Its 17 features describe shape (area, bounding box, elongation,
solidity), first-stage probability (mean, maximum, 90th percentile, mean in a 20-pixel ring),
darkness (inside, ring, contrast and spread, relative to the disk median), the radial position
on the disk and the number of candidates in the image. It learns from the 5,081 out-of-fold
refined candidates. The features use the mean map of one fold's 1024 and 2048 models, as in
training; at test time each candidate is scored with each fold's map and the five scores are
averaged. Candidates with a score below 0.35 are removed.

\begin{figure*}[t]
\centering
\includegraphics[width=\textwidth]{figures/examples.png}
\caption{Out-of-fold example segmentations with low, median and high per-image PQ. Annotators
in cyan, yellow and magenta; our masks in red. Left: the annotator marks a long faint filament
that we miss, and our masks cover other parts of the marked filaments than the annotator's
outlines. Middle: three annotators and our masks agree on three filaments; one annotator also
marks a small filament at the top that the others and our model leave out. Right: close
agreement.}
\label{fig:examples}
\end{figure*}

% ==================================== %
%              EVALUATION              %
% ==================================== %
\section{Evaluation}\label{sec:evaluation}

\subsection{Strategy}

We judge every change by out-of-fold PQ with the fixed post-processing. Two training seeds of
the same model differ by up to 0.004 on one fold, so we kept a change only if it gained more
than 0.007 on two folds. The filter is judged by nested cross-validation: for each outer
fold, it is trained on the other four folds, its threshold is chosen by an inner 4-fold
cross-validation on those folds, and both are applied to the outer fold, so no score uses a
threshold tuned on its own images. For the final pipeline we report per-fold scores, the
spread over folds, and a bootstrap over images (2,000 samples) for pooled PQ and for the
paired gains. The organisers' own plotting code draws the distributions and matching counts.

\subsection{Results}

\begin{table}[t]
\caption{Out-of-fold PQ per fold, mean and standard deviation over folds.}
\label{tab:folds}
\centering
\small
\begin{tabular}{lccccc|c}
\toprule
Pipeline & F0 & F1 & F2 & F3 & F4 & mean $\pm$ sd \\
\midrule
1024 model & .410 & .421 & .406 & .406 & .422 & .413 $\pm$ .008 \\
2048 model & .415 & .431 & .422 & .423 & .431 & .425 $\pm$ .007 \\
1024 + 2048 & .420 & .434 & .419 & .423 & .436 & .426 $\pm$ .008 \\
+ seed refiner & .427 & .440 & .419 & .428 & .443 & .432 $\pm$ .010 \\
+ learned filter & .446 & .446 & .435 & .442 & .456 & .445 $\pm$ .007 \\
\bottomrule
\end{tabular}
\end{table}

\begin{figure}[t]
\centering
\includegraphics[width=0.92\columnwidth]{figures/folds_errorbars.png}
\caption{Out-of-fold PQ per pipeline: mean $\pm$ standard deviation over the five folds
(blue) and each fold (gray).}
\label{fig:folds}
\end{figure}

Table~\ref{tab:folds} and Fig.~\ref{fig:folds} give the scores. The fold spread (0.007--0.010)
is larger than the gains, because folds differ in difficulty; the paired comparisons are
clearer. The 1024 + 2048 mean beats the 1024 model on every fold (+0.013 on average). The
refiner gains on four folds and is neutral on fold 2; the filter gains on all five folds
(+0.006 to +0.018). Pooled over all 707 images, PQ is 0.426 [0.416, 0.435] for the first stage,
0.431 [0.422, 0.440] after the refiner and 0.445 [0.435, 0.454] after the filter (95\% bootstrap
intervals). The paired gains are +0.0054 [0.0032, 0.0078] for the refiner and +0.0136 [0.0106,
0.0168] for the filter, both positive in all 2,000 samples. These intervals cover the choice of
images, not training noise (about 0.003 between refiner seeds). On the public leaderboard the
four pipelines scored 0.35, 0.36, 0.37 and 0.38; the public score never moved against the local
score and stays about 0.06 below it. On the 296 images with two or three annotators, our masks
score PQ 0.446 (SQ 0.677) against each annotator, while one annotator scores 0.340 (SQ 0.634)
against the others: the model agrees with each annotator more than the annotators agree with
each other.

\begin{figure}[t]
\centering
\includegraphics[width=\columnwidth]{figures/iou_distribution.png}
\includegraphics[width=\columnwidth]{figures/matching_counts.png}
\caption{Out of fold, all 707 images, drawn with the organisers' plotting code. Top: IoU of
overlapping (GT, prediction) pairs. Bottom: matching cases; ``2:1'' is a prediction that
overlaps two GT filaments, ``1:2'' a GT filament that overlaps two predictions.}
\label{fig:official}
\end{figure}

Fig.~\ref{fig:official} compares the first stage with the final pipeline. Refiner and filter
together remove 1,411 false positives (3,437 $\to$ 2,026) at the cost of 350 more false
negatives (2,870 $\to$ 3,220). The GT filaments that overlap two or more predictions fall from
394 to 242; the predictions that overlap two or more GT filaments stay about the same (257
$\to$ 262; the refiner alone raised them to 281). The mean IoU of overlapping pairs rises from
0.570 to 0.590 (Dice 0.703 $\to$ 0.722).

\subsection{Strengths, weaknesses and possible improvements}
\label{sec:weak}

\textbf{Strengths.} The pipeline is simple, has few tuned values, and every choice was
validated out of fold on two or more folds. It runs at 1.5\,s per image on one RTX~3070
(0.83\,s for the 2048 models) and 79\,s on four CPU threads. On a GPU it is deterministic: the
notebook in our repository rebuilds the submitted CSV with the same SHA-256 from the stored
weights on the same GPU type; on another GPU type the masks agree with a mean IoU of 0.9998,
and on a CPU with 0.997.

\textbf{Weaknesses: where the remaining PQ goes.} Table~\ref{tab:gap} splits the errors of the
final pipeline (out of fold, all 707 images; \texttt{scripts/gap\_analysis.py}). Of the 1,878
GT filaments that no prediction touches, 683 have no first-stage pixel above 0.5, 574 have
fewer than 300 such pixels, and 271 were found by the refiner but removed by the filter, which
pays for 1,000 fewer false positives. On images with two or three annotators, only 56\% of the
missed filaments are also marked by another annotator, against 93\% of the matched ones: about
half of the misses are one annotator's choice (Fig.~\ref{fig:disagree}, middle). Near misses go
both ways: the prediction is smaller than the GT in 462 cases and larger in 370; 129 GT
filaments are split and 189 predictions span two GT filaments (Fig.~\ref{fig:disagree},
right). Growing or shrinking every mask by one pixel lowers PQ to 0.418 and 0.417, so there is
no global bias to correct.

\begin{table}[t]
\caption{Error types of the final pipeline (out of fold) and the PQ if every case of the type
were fixed (missed GT found at IoU 0.6, near misses at IoU 0.55, false positives removed).}
\label{tab:gap}
\centering
\small
\begin{tabular}{lrc}
\toprule
Error type & count & PQ if fixed \\
\midrule
Final pipeline (TP) & 4,979 & .445 \\
GT never touched (IoU $<$ 0.05) & 1,878 & .528 \\
Near-miss pairs (IoU 0.3--0.5) & 916 & .510 \\
GT with weak overlap (0.05--0.3) & 426 & -- \\
FP with no GT nearby & 697 & .466 \\
\bottomrule
\end{tabular}
\end{table}

\begin{figure*}[t]
\centering
\includegraphics[width=\textwidth]{figures/disagreement.png}
\caption{Zoomed out-of-fold cases from images with three annotators (cyan, yellow, magenta; our
masks in red), each the largest GT filament of its kind. Left: the annotators' outlines of one
filament differ widely (best agreement IoU 0.31) and our mask lies between them. Middle: one
annotator marks a long faint filament near the limb that the other two and our model leave
out. Right: a near miss where the annotators agree (IoU 0.72): they split the filament from its
lower arm, our mask joins them (IoU 0.47).}
\label{fig:disagree}
\end{figure*}

\textbf{What did not help} (all out of fold, same post-processing): a majority-vote or
random-annotator target instead of the soft mean ($-0.003$, $-0.005$); a mean-probability
filter on candidates (+0.003); eight flips and rotations as test-time augmentation (+0.004 and
$-0.003$ on two folds, at four times the cost); a ResNet50 encoder at 1024 (+0.004 as a third
model, +0.006 instead of the ResNet34); averaging two refiner seeds (no gain); a focal + Dice
loss ($-0.005$); six more filter features (two per-resolution probabilities, their
difference, edge strength, ring spread, dark quantile: $-0.0004$); an auxiliary spine output trained on the MAGFiLO spines (+0.005 and +0.003
for the 1024 model alone, but only +0.001 inside the 1024 + 2048 mean); more candidates for the
filter (first-stage threshold 0.4, minimum area 150: $\pm$0.000 pooled); input cleaning as in
EdgeAttNet~\cite{edgeattnet} (disk mask, radial flattening, CLAHE: $-0.004$ and $-0.002$); an
auxiliary output for the contact zone between filaments (the zone is present in only 121 of
707 images and the output stays at zero); and other weights than 0.5 for the 2048 map (best
0.5).

\textbf{Possible improvements.} A refiner that uses the spine to cover the whole filament
(spine supervision of the first stage did not help); an instance signal that reaches further
than a contact zone, such as a per-filament embedding, to split masks that span two filaments;
and self-supervised pretraining on unlabeled H$\alpha$ archives, which
the organisers allow when it excludes the test observations.

% ==================================== %
%            ACKNOWLEDGMENT            %
% ==================================== %
\section{Acknowledgment}\label{sec:acknowledgment}
    % Do not change
    The organization of the Filament Segmentation Challenge 2026 is partially supported by the U.S. National Science Foundation (NSF) under Grant No. 2209912 and 2433781, directorate for Computer and Information Science and Engineering (CSE), and Office of Advanced Cyberinfrastructure (OAC), and Grant No. 2511630, AST Division Of Astronomical Sciences and MPS Directorate for Mathematical and Physical Sciences.

    The Filament Segmentation Challenge 2026 is sponsored by the U.S. National Science Foundation (NSF) National Solar Observatory (NSO). 

    This work utilizes GONG data obtained by the NSO Integrated Synoptic Program, managed by the National Solar Observatory, which is operated by the Association of Universities for Research in Astronomy (AURA), Inc. under a cooperative agreement with the National Science Foundation and with contribution from the National Oceanic and Atmospheric Administration. The GONG network of instruments is hosted by the Big Bear Solar Observatory, High Altitude Observatory, Learmonth Solar Observatory, Udaipur Solar Observatory, Instituto de Astrofísica de Canarias, and Cerro Tololo Interamerican Observatory.

\bibliographystyle{ACM-Reference-Format}
\bibliography{main}

\end{document}
